\section{Mechanism counterexamples and their sufficient conditions}
\subsection{An unexcited storage clamp: E2}
Two ideal worlds receive 20~$\mu$W for 2 s and face a 30~$\mu$J service
transaction at 2.5 s. Both start with zero stored energy, use
$C=10~\mu$F and must retain the energy associated with a 1 V minimum rail.
Their clamps are 3 V and 5 V. The current maximum voltage is
$\sqrt{8}=2.828$ V, so neither clamp activates. Their complete current
observed event and rail traces coincide and both services pass.

The common legal change sets $C=4~\mu$F within the declared
$[4,12]~\mu$F domain. At 3 V, only 18~$\mu$J can be stored; 22~$\mu$J
spills and only 16~$\mu$J is available above the 1 V reserve. Service
fails with a $-14~\mu$J margin. At 5 V, all 40~$\mu$J is retained and
the margin is $8~\mu$J, so service passes. The task, source, requested
transaction and common change are identical. Thus the current trace does
not distinguish a mechanism needed for the prospective verdict~\cite{structural}.

An enriched static description resolves the ambiguity. Under the exact
assumptions that all input precedes service, the capacitor is ideal and
there is no leakage or intervening workload, define
\begin{equation}
 M(C)=\min\!\left(E_0+P_{\rm in}t_{\rm in},\frac{CV_{\max}^2}{2}\right)
       -\frac{CV_{\min}^2}{2}-E_{\rm service}.
 \label{eq:margin}
\end{equation}
Service passes exactly when $M(C)\geq0$. The first term is the stored
energy after the source ends; the next is the required residual energy.
With no other transaction between those times, their difference is exactly
the budget available to service. The sufficient object is the functional
$M(\cdot)$ and its parameters on the legal domain, not merely its value at
the current $C$. Both worlds have the same current margin. A current scalar
margin alone therefore retains the counterexample.

\subsection{An unexcited reset transition: E3}
The second pair captures event 1 at 1 s and attempts its single transmission
at 4 s, before a 5 s deadline. A volatile pending memory and a retained
pending memory both deliver in the current uninterrupted execution. The
same legal change inserts a power outage from 2 to 3 s. Volatile memory
clears the pending event and delivers none; retained memory delivers one.
External event timing and the transmission schedule remain fixed.
This witness changes a state transition rather than scaling an energy
total~\cite{structural}.

The sufficient static state description must state capture, reset and send
semantics and their order. If pending state survives each intervening reset,
and the scheduled send can execute after restoration, the captured event
remains available. If the only pending copy is cleared and no recapture is
possible before sending, it does not. These are conditions of the ideal
machine, not a claim that an arbitrary nonvolatile device has free writes,
atomic persistence or a qualified energy budget.

% Generated from hash-verified saved data; no scientific execution.
\begin{table}[htbp]
\centering\small
\caption{E2--E3 structural witnesses: both worlds pass before a common legal intervention, then separate. These are two exposed ideal pairs, not statistical repetitions.}
\label{tab:witnesses}
\begin{tabular}{@{}llccr@{}}\toprule
Mechanism & World & Current & Changed & Changed events \\\midrule
Clamp & low-clamp & PASS & FAIL & 0 \\
Clamp & high-clamp & PASS & PASS & 1 \\
Reset memory & volatile & PASS & FAIL & 0 \\
Reset memory & retained & PASS & PASS & 1 \\
\bottomrule\end{tabular}
\end{table}


The saved analytic and numerical clamp implementations and the separate
memory interpreter agree on eight current/intervention cases at both 1 ms
and 0.5 ms settings: sixteen resolution-specific comparisons. Maximum
reported energy discrepancy is $2.84\times10^{-12}~\mu$J. The enriched
static model agrees on all 28 exposed grid points, and an additive static
control agrees on five service costs spanning PASS and FAIL. Removing the
unexcited clamp or reset-retention fact restores aliasing in these encodings.
These checks establish software consistency under common ideal assumptions,
not independent physical truth or universal necessity of a particular field
name~\cite{structural}.

\subsection{Connecting the witnesses to transient BF behavior}
Seven separate saved BF controls exercise source timing, startup, loaded
rail, reset and reception (Table~\ref{tab:reference}). Their transient model
uses ready/off thresholds, boot energy, finite TX current and duration, and
the loaded-rail relation
\begin{equation}
 V_{\rm load}=\sqrt{2E/C}-I_{\rm TX}R_{\rm ESR}.
\end{equation}
A failed energy or rail gate can reset the program and clear volatile
pending state. The causal path is load to stored energy and loaded rail,
then reset and program behavior. The source is prescribed, so no
load-dependent source feedback is established~\cite{mechanisms}.

% Generated from hash-verified saved data; no scientific execution.
\begin{table}[htbp]
\centering\small
\caption{Seven saved digital BF reference controls. First reception after 100 ms fails the deadline; a passing rail gate without a received event is insufficient. Physical qualification is UNKNOWN throughout.}
\label{tab:reference}
\begin{tabular}{@{}lrrcc@{}}\toprule
Control & Input ($\mu$J) & RX (ms) & Service & Rail \\\midrule
Fast pulse & 150 & 15.8 & PASS & PASS \\
Slow pulse & 150 & 54.8 & PASS & PASS \\
High ESR & 150 & -- & FAIL & FAIL \\
No energy & 0 & -- & FAIL & PASS \\
Late pulse & 150 & 105.8 & FAIL & PASS \\
Latched reset & 210 & 46.9 & PASS & PASS \\
Volatile reset & 210 & -- & FAIL & PASS \\
\bottomrule\end{tabular}
\end{table}


Fast, slow and late pulses each supply 150~$\mu$J, yet their first
receptions differ; the late one misses the 100 ms deadline. A 150~$\Omega$
ESR control prevents reception despite the same source integral. The
no-energy control passes its rail check vacuously and fails service.
The reset pair supplies 210~$\mu$J with the same 40~$\mu$J shunt and
postponed transmission, but pending-event survival differs. Its retained
mechanism is an explicitly ideal passive mechanical latch able to capture
without stored electrical energy. It is not a model of a qualified
commercial contact device or a powered NVM write.

All seven agree with their predeclared control expectations, including the
intended failures. That agreement is different from seven service passes.
The full input laws and source calibration limits accompany the saved data;
physical qualification remains UNKNOWN for every control. These transient
examples connect the structural issue to BF operation without turning the
two ideal witness models into a complete physical node.
\FloatBarrier
